% SPDX-License-Identifier: PMPL-1.0-or-later
% Copyright (c) 2026 Jonathan D.A. Jewell (hyperpolymath) <j.d.a.jewell@open.ac.uk>
%
% Towards a Decalogue of Database Query Safety:
% Progressive Type-Level Guarantees from Parse-Time to Provenance
%
% Draft manuscript under audit

\documentclass[11pt,a4paper]{article}

% ============================================================================
% Packages
% ============================================================================
\usepackage[utf8]{inputenc}
\usepackage[T1]{fontenc}
\usepackage{lmodern}
\usepackage{amsmath,amssymb,amsthm}
\usepackage{mathpartir}
\usepackage{stmaryrd}
\usepackage{listings}
\usepackage{xcolor}
\usepackage{hyperref}
\usepackage{cleveref}
\usepackage{booktabs}
\usepackage{multirow}
\usepackage{graphicx}
\usepackage{tikz}
\usetikzlibrary{arrows.meta,positioning,fit,calc,decorations.pathmorphing}
\usepackage{enumitem}
\usepackage{subcaption}
\usepackage{balance}
\usepackage[margin=1in]{geometry}
\usepackage{doi}
\usepackage{natbib}
\bibliographystyle{plainnat}

% ============================================================================
% Theorem environments
% ============================================================================
\newtheorem{theorem}{Theorem}[section]
\newtheorem{lemma}[theorem]{Lemma}
\newtheorem{proposition}[theorem]{Proposition}
\newtheorem{corollary}[theorem]{Corollary}
\newtheorem{definition}[theorem]{Definition}
\newtheorem{example}[theorem]{Example}
\newtheorem{remark}[theorem]{Remark}

% ============================================================================
% Listings configuration
% ============================================================================
\definecolor{kwcolor}{RGB}{0,80,160}
\definecolor{strcolor}{RGB}{160,32,32}
\definecolor{commentcolor}{RGB}{80,120,80}
\definecolor{typecolor}{RGB}{128,0,128}
\definecolor{bglight}{RGB}{248,248,248}

\lstdefinelanguage{VQL}{
  keywords={SELECT,FROM,WHERE,TYPED,BY,RETURNS,EXPECTING,EFFECTS,FRESH,
            WITHIN,CONSUME,AFTER,USE,SESSION,USAGE,LIMIT,PROOF,
            EXACTLY,AT_MOST,AT_LEAST,MANY,BETWEEN,AND,OR,NOT,
            HEXAD,FEDERATION,STORE,GRAPH,VECTOR,TENSOR,SEMANTIC,
            DOCUMENT,TEMPORAL,PROVENANCE,SPATIAL,AS,OF,INSERT,
            UPDATE,DELETE,SET,ORDER,GROUP,HAVING,ASC,DESC,
            FETCH,UNIQUE,OPTIONAL,COALESCE,OPEN,CLOSE,BEGIN,
            COMMIT,ROLLBACK,TRANSACTION,CONNECTION,USING,REMOVE,
            WITH,IN,IS,NULL,NOW,SECONDS},
  keywordstyle=\color{kwcolor}\bfseries,
  ndkeywords={Read,Write,DDL,Network,Proof,Index,Int,Float,String,
              Bool,UUID,Timestamp,Duration,Bytes,Array,Set,Map,
              ExactlyOne,AtMostOne,NonEmpty,List,Nullable,Temporal,
              Freshness,Session,Connection},
  ndkeywordstyle=\color{typecolor}\bfseries,
  comment=[l]{--},
  morecomment=[s]{/*}{*/},
  commentstyle=\color{commentcolor}\itshape,
  stringstyle=\color{strcolor},
  string=[b]",
  basicstyle=\small\ttfamily,
  breaklines=true,
  showstringspaces=false,
  tabsize=2,
  backgroundcolor=\color{bglight},
  frame=single,
  framerule=0.4pt,
  rulecolor=\color{gray!40},
}

\lstdefinelanguage{Idris2}{
  keywords={data,where,total,public,export,module,import,let,in,do,
            case,of,if,then,else,the,rewrite,with,impossible,
            interface,implementation,using,record,namespace,
            mutual,parameters,partial,covering,MkTemporal},
  keywordstyle=\color{kwcolor}\bfseries,
  ndkeywords={Type,Nat,Vect,DPair,Maybe,Just,Nothing,IO,Unit,
              Schema,Table,Column,Query,Result,Cardinality,
              Effect,TemporalBound,SessionState,Connection,
              Ready,InTransaction,Closed,Read,Write},
  ndkeywordstyle=\color{typecolor}\bfseries,
  comment=[l]{--},
  morecomment=[s]{\{-}{-\}},
  commentstyle=\color{commentcolor}\itshape,
  stringstyle=\color{strcolor},
  string=[b]",
  basicstyle=\small\ttfamily,
  breaklines=true,
  showstringspaces=false,
  tabsize=2,
  backgroundcolor=\color{bglight},
  frame=single,
  framerule=0.4pt,
  rulecolor=\color{gray!40},
  literate={->}{{$\to$}}2 {=>}{{$\Rightarrow$}}2,
}

\lstset{
  basicstyle=\small\ttfamily,
  breaklines=true,
  showstringspaces=false,
}

% ============================================================================
% Macros
% ============================================================================
\newcommand{\VQL}{\textsc{VQL-UT}}
\newcommand{\VeriSimDB}{\textsc{VeriSimDB}}
\newcommand{\TypeLL}{\textsc{TypeLL}}
\newcommand{\PanLL}{\textsc{PanLL}}
\newcommand{\LevelN}[1]{\textsf{L#1}}
\newcommand{\Safe}[1]{\ensuremath{\mathsf{Safe}_{#1}}}
\newcommand{\Leq}{\ensuremath{\sqsubseteq}}
\newcommand{\Lat}{\ensuremath{\mathcal{L}}}
\newcommand{\Sem}[1]{\ensuremath{\llbracket #1 \rrbracket}}

% ============================================================================
% Title
% ============================================================================
\title{%
  \textbf{Towards a Decalogue of Database Query Safety:\\
  Progressive Type-Level Guarantees\\
  from Parse-Time to Provenance}%
}

\author{%
  Jonathan D.A. Jewell\\
  \texttt{j.d.a.jewell@open.ac.uk}\\
  \medskip
  \small{Independent Researcher}
}

\date{March 2026}

% ============================================================================
% Document
% ============================================================================
\begin{document}

\maketitle

% ============================================================================
% Abstract
% ============================================================================
\begin{abstract}
Database query bugs span a remarkably wide spectrum, from syntactic
malformation to cryptographic provenance failure.  We observe that these
failure modes form a \emph{total order} by severity and subtlety: each
successive class of bug presupposes the absence of all preceding classes.
We develop this observation into a \emph{draft decalogue}---ten progressive
levels of type-level query safety intended to subsume the prior ones.

Levels 1--6 (parse-time safety, schema-binding safety, type-compatible
operations, null safety, injection-proof safety, and permission safety)
are individually addressed, though rarely unified, in the existing
literature.  Levels 7--10 (cardinality safety, consistency safety,
temporal safety, and provenance safety) are proposed here as candidate
extensions to that safety envelope.  We give formal definitions,
an intended subsumption structure, and a prototype implementation plan
for representing them in a dependently-typed setting.

We present \VQL{}, a query language for the multi-modal database
\VeriSimDB{}, as a design target and early prototype rather than a finished
system in which all ten levels are already enforced at compile time.
The current repository contains Rust formatter/linter tooling, draft Idris\,2
encodings for parts of the safety model and ABI, and an implementation plan
for the remaining pipeline.  The paper's strongest claims are therefore about
the taxonomy, formal problem framing, and proposed architecture; full
mechanisation and end-to-end enforcement remain ongoing work.
\end{abstract}

\noindent
\textbf{Status note:} Draft manuscript under repository audit. Earlier drafts
overstated implementation maturity; this version should be read as a proposal
and prototype report, not as a claim that the repository already provides a
fully verified end-to-end VQL-UT system.

\noindent
\textbf{Keywords:} query safety, dependent types, type-safe databases,
multi-modal databases, formal verification, provenance, temporal queries

\medskip
\noindent
\textbf{ACM CCS:} Software and its engineering $\to$ Formal language
definitions; Information systems $\to$ Query languages; Security and
privacy $\to$ Database and storage security

% ============================================================================
\section{Introduction}\label{sec:intro}
% ============================================================================

The history of database security and correctness is, in large part, a
history of classification failure.  SQL injection
\citep{halfond2006classification}, null-related errors
\citep{hoare2009null}, and schema drift
\citep{curino2008automating} are studied in isolation, addressed by
different tools, and defended by different communities.  Yet these
failure modes are not independent: an injection attack presupposes
that the attacker's input survives parsing (Level~1) and is bound to a
schema context (Level~2).  A cardinality violation presupposes that the
query is well-typed (Level~3) and null-safe (Level~4).  The failures
form a chain, not a cloud.

This paper makes three contributions:

\begin{enumerate}[label=(\roman*)]
  \item \textbf{The Decalogue.}  We identify ten levels of database
    query safety, numbered \LevelN{1} through \LevelN{10}, and show
    how they are intended to form a total order under the subsumption relation: a
    query that is \Safe{k}-safe is automatically \Safe{j}-safe for all
    $j < k$.

  \item \textbf{The Novel Quartet.}  While levels 1--6 are
    individually addressed (though seldom unified) in the literature,
    levels 7--10---cardinality safety, consistency safety, temporal
    safety, and provenance safety---have never been formalised as
    type-level invariants in a query language.  We provide draft formal
    definitions and argue that they can be encoded in a dependently-typed
    core calculus.

  \item \textbf{A Concrete Realisation.}  We present \VQL{}, a query
    language for the \VeriSimDB{} multi-modal database system, in which
    the long-term goal is compile-time enforcement across all ten levels.
    The current repository provides an early prototype: Rust surface
    tooling, draft Idris\,2 encodings, and ABI/FFI scaffolding rather
    than a finished end-to-end verified implementation.
\end{enumerate}

\paragraph{Organisation.}
\Cref{sec:background} surveys related work.
\Cref{sec:decalogue} presents the ten levels with formal definitions.
\Cref{sec:lattice} establishes the cumulative safety ordering.
\Cref{sec:typesystem} describes the type system design.
\Cref{sec:multimodal} discusses integration with multi-modal databases.
\Cref{sec:implementation} covers the implementation architecture.
\Cref{sec:comparison} compares \VQL{} with existing query languages.
\Cref{sec:future} outlines future work.
\Cref{sec:conclusion} concludes.

% ============================================================================
\section{Background and Related Work}\label{sec:background}
% ============================================================================

\subsection{SQL Injection and Prepared Statements}

SQL injection was first publicly described by Rain Forest Puppy in
Phrack Magazine in 1998 \citep{rfp1998nt} and has since become the
single most prevalent class of web application vulnerability, holding
the \#1 position in the OWASP Top Ten for over a decade
\citep{owasp2021top10}.  The Equifax breach of 2017 exposed the
personal data of 147 million individuals through an injection vector
\citep{usgao2018equifax}.

Prepared statements, introduced as a mitigation in the late 1990s,
separate query structure from data by using parameter placeholders.
While effective against first-order injection, they do not address
second-order injection (where previously stored data is interpolated
into a new query) and rely on correct usage by the programmer---a form
of \emph{safety by convention} rather than \emph{safety by
construction}.

\VQL{} addresses this at Level~5 (injection-proof safety) by making
the separation structural: user input and query syntax inhabit
different types that cannot be combined.  There is no function
$\mathit{String} \to \mathit{Query}$; the only path from text to query
passes through the parser, which enforces all lower levels.

\subsection{ORM Safety and Query Builders}

Object-relational mappers (ORMs) such as Hibernate \citep{bauer2015java},
Django ORM \citep{holovaty2009django}, and Prisma
\citep{prisma2023docs} provide varying degrees of type safety.
Prisma, in particular, generates TypeScript types from the database
schema, achieving something close to our Level~2 (schema-binding) and
Level~6 (permission/result-type safety).  However, ORMs typically
operate at the application language level and cannot enforce
guarantees \emph{within} the query language itself.

Slick \citep{slick2023docs} and Quill \citep{ioffe2017quill}
take a more principled approach by embedding queries into Scala's type
system, achieving Levels~1--3 and partial Level~4.  jOOQ
\citep{jooq2023docs} generates Java code from database schemas,
providing strong schema-binding but relying on Java's non-dependent
type system, which limits expressiveness at higher levels.

\subsection{Dependent Types for Databases}

The application of dependent types to database query safety has been
explored in several contexts.  \citet{oury2008power} demonstrated
dependent types for data formats in Epigram.
\citet{brady2013idris} showed how Idris's dependent types can encode
protocol correctness, a technique we extend to database session types
at Level~10.  \citet{chlipala2010ur} presented Ur/Web, which uses a
type system with row types and type-level computation to verify web
applications against database schemas, achieving approximately our
Levels~1--6.

More recently, \citet{weirich2017dependent} surveyed dependent types in
Haskell, and \citet{cockx2021idris2} described Idris\,2's quantitative
type theory (QTT), which we exploit for linear resource tracking at
Level~10.

The key distinction of our work is \emph{unification}: where prior
systems address subsets of the safety levels, \VQL{} addresses all ten
within a single type-theoretic framework, and crucially, formalises the
subsumption ordering that relates them.

\subsection{Type-Safe Query Languages}

LINQ \citep{meijer2006linq} introduced language-integrated queries to
C\#, providing compile-time type checking of query expressions.
EdgeQL \citep{edgedb2023docs} embeds queries deeply into EdgeDB's
type system.  Malloy \citep{malloy2023docs} provides a semantic
modelling layer with compile-time type checking.  PRQL
\citep{prql2023docs} focuses on ergonomic pipelined query syntax with
type inference.

GraphQL \citep{graphql2023spec} provides schema-aware query validation
and result typing (Levels~1--3 and~6), but its type system lacks null
safety (field nullability is a common source of client errors) and has
no notion of effects, cardinality, temporality, or linearity.

None of these systems formalise a progressive safety ordering, and
none address Levels~7--10.

\subsection{Provenance and Lineage}

Data provenance---tracking the origin, transformation, and custody
of query results---has been studied extensively in the database
community \citep{cheney2009provenance,buneman2001and,green2007provenance}.
The semiring framework of \citet{green2007provenance} provides an
algebraic foundation for provenance in relational queries.  However,
provenance is typically treated as a runtime annotation rather than a
type-level invariant.

Our Level~10 (provenance safety) encodes provenance as a type-level
property: query results carry cryptographic attestations of their
origin, and the type system ensures that provenance chains are
maintained through compositions.

% ============================================================================
\section{The Decalogue: Ten Levels of Query Safety}\label{sec:decalogue}
% ============================================================================

We now present the ten safety levels.  Each level is defined by a
\emph{property} (what is guaranteed), a \emph{mechanism} (how the
guarantee is enforced), and a \emph{failure mode} (what goes wrong
without it).  Each level \emph{strictly subsumes} all preceding
levels: a Level~$k$ guarantee is only meaningful if Levels~1 through
$k-1$ are already established.

\begin{definition}[Safety Level]
A \emph{safety level} $\ell \in \{1, \ldots, 10\}$ is a predicate
$\Safe{\ell}(q)$ on queries $q$ such that:
\[
  \Safe{k}(q) \implies \Safe{j}(q) \quad \text{for all } 1 \leq j < k \leq 10
\]
\end{definition}

\noindent
We write $\LevelN{k}$ to refer to safety level $k$ and $\Safe{k}$ to
refer to the predicate that characterises it.

% ----------------------------------------------------------------------------
\subsection{Level 1: Parse-Time Safety}\label{sec:L1}
% ----------------------------------------------------------------------------

\begin{definition}[\Safe{1}: Parse-Time Safety]
A query $q$ is \Safe{1}-safe if and only if $q$ is a member of the
language $\mathcal{L}(G)$ defined by the VQL-UT grammar $G$.
Formally:
\[
  \Safe{1}(q) \iff q \in \mathcal{L}(G)
\]
where $G$ is a context-free grammar and $\mathcal{L}(G)$ its generated
language.
\end{definition}

Parse-time safety is the foundation.  It ensures that only
syntactically well-formed queries enter the processing pipeline.
The parser is the first line of defence: it rejects malformed input
before any semantic analysis occurs.

\paragraph{Mechanism.}
The \VQL{} parser, written in ReScript, implements a recursive-descent
parser for the VQL-UT grammar (defined in EBNF).  A successful parse
produces a typed abstract syntax tree (AST); failure produces a
structured error with position information and suggestions.

\paragraph{Failure mode without \LevelN{1}.}
Raw strings reach the database engine.  The engine must either reject
them (losing the opportunity for helpful error messages) or attempt
ad-hoc parsing (introducing undefined behaviour on malformed input).
The canonical example is Bobby Tables \citep{xkcd327}: a student's
name containing raw SQL causes a school database to execute
\texttt{DROP TABLE students}.

\begin{example}
\begin{lstlisting}[language=VQL]
-- Rejected at Level 1: syntax error
SELCT users WHERE id = 1
-- Error: Unknown keyword "SELCT" at position 0.

-- Accepted at Level 1: syntactically valid
SELECT DOCUMENT.name FROM HEXAD 3f2504e0-... WHERE DOCUMENT.id = $id :: Int
\end{lstlisting}
\end{example}

% ----------------------------------------------------------------------------
\subsection{Level 2: Schema-Binding Safety}\label{sec:L2}
% ----------------------------------------------------------------------------

\begin{definition}[\Safe{2}: Schema-Binding Safety]
A query $q$ is \Safe{2}-safe if and only if $q$ is \Safe{1}-safe and
every identifier in $q$ that references a schema element (table,
column, modality, relationship) resolves to an entity in the current
schema $\Sigma$:
\[
  \Safe{2}(q) \iff \Safe{1}(q) \land \forall\, r \in \mathrm{refs}(q).\; r \in \mathrm{dom}(\Sigma)
\]
where $\mathrm{refs}(q)$ is the set of schema references in $q$ and
$\mathrm{dom}(\Sigma)$ is the domain of the schema.
\end{definition}

\paragraph{Mechanism.}
The schema binder, implemented in Idris\,2, represents the database
schema as a dependent type.  The type itself encodes which tables,
columns, and modalities exist:

\begin{lstlisting}[language=Idris2]
data Schema : Type where
  MkSchema : (tables : Vect n (Table cols)) -> Schema

data Table : Vect m Column -> Type where
  MkTable : (name : String) -> (columns : Vect m Column) -> Table columns
\end{lstlisting}

A query that references a non-existent column fails type checking,
because the dependent type requires a proof that the column exists
in the table's column vector.

\paragraph{Failure mode without \LevelN{2}.}
Queries reference columns that do not exist and silently return empty
results.  A typo in \texttt{order\_totl} instead of
\texttt{order\_total} produces zero rows with no error---a failure
mode that can persist undetected for months in production
\citep{kim2013silent}.

\begin{example}
\begin{lstlisting}[language=VQL]
-- Schema: users(id: UUID, name: String, email: String?)
SELECT DOCUMENT.emial FROM HEXAD ...
  TYPED BY users@1.0
-- Level 2 error: "emial" is not a column in "users".
-- Available: id, name, email.
\end{lstlisting}
\end{example}

% ----------------------------------------------------------------------------
\subsection{Level 3: Type-Compatible Operations}\label{sec:L3}
% ----------------------------------------------------------------------------

\begin{definition}[\Safe{3}: Type-Compatible Operations]
A query $q$ is \Safe{3}-safe if and only if $q$ is \Safe{2}-safe and
every operation $\mathit{op}(e_1, \ldots, e_n)$ in $q$ satisfies the
type compatibility constraint:
\[
  \Safe{3}(q) \iff \Safe{2}(q) \land
  \forall\, \mathit{op}(e_1, \ldots, e_n) \in q.\;
  (\tau(e_1), \ldots, \tau(e_n)) \in \mathrm{dom}(\mathit{op})
\]
where $\tau(e)$ is the type of expression $e$ and $\mathrm{dom}(\mathit{op})$
is the set of valid argument type tuples for operator $\mathit{op}$.
\end{definition}

\paragraph{Mechanism.}
The type checker (Idris\,2) verifies that every comparison, arithmetic
operation, aggregation, and function application receives arguments
of compatible types.  Type compatibility is defined by a decidable
relation, and the checker is \emph{total}: it always terminates with
a definitive accept-or-reject verdict.

\paragraph{Failure mode without \LevelN{3}.}
Implicit type coercion produces semantically meaningless results.  The
expression \texttt{date\_of\_birth > 42} might compare a date to an
integer via implicit string conversion, yielding results that are
technically consistent with the coercion rules but violate every
reasonable expectation.  The Therac-25 radiation therapy incidents
(1985--1987), in which mode flags and dosage counters were conflated
due to absent type distinctions, resulted in patient deaths
\citep{leveson1993therac}.

\begin{example}
\begin{lstlisting}[language=VQL]
-- Rejected at Level 3: type incompatibility
SELECT DOCUMENT.name FROM HEXAD ...
  WHERE DOCUMENT.name > 42
-- Error: operator (>) applied to (String, Int). Expected compatible types.

-- Accepted at Level 3
SELECT DOCUMENT.age FROM HEXAD ...
  WHERE DOCUMENT.age :: Int > 42
\end{lstlisting}
\end{example}

% ----------------------------------------------------------------------------
\subsection{Level 4: Null Safety}\label{sec:L4}
% ----------------------------------------------------------------------------

\begin{definition}[\Safe{4}: Null Safety]
A query $q$ is \Safe{4}-safe if and only if $q$ is \Safe{3}-safe and
every expression $e$ in $q$ with type $\mathsf{Nullable}(\tau)$ is
either (a)~explicitly handled by a null-eliminating operation
(\texttt{COALESCE}, \texttt{IS NOT NULL} narrowing, or pattern
matching), or (b)~propagated with the nullable type intact:
\[
  \Safe{4}(q) \iff \Safe{3}(q) \land \forall\, e \in q.\;
  \mathit{nullable}(e) \implies \mathit{handled}(e)
\]
\end{definition}

\paragraph{Mechanism.}
The type system tracks nullability as a type-level wrapper.  A column
declared \texttt{Nullable(Text)} is a fundamentally different type from
\texttt{Text}.  Operations that require non-null input reject nullable
arguments at compile time.  \texttt{LEFT JOIN} automatically lifts all
columns from the right-hand table to their nullable counterparts.

The nullability algebra follows three-valued logic faithfully: if $e$
has type $\mathsf{Nullable}(\tau)$, then $e + 1$ has type
$\mathsf{Nullable}(\tau)$ (null propagation), and $e = e$ has type
$\mathsf{Nullable}(\mathsf{Bool})$ (SQL's null~$=$~null evaluates to
null, not true).

\paragraph{Failure mode without \LevelN{4}.}
Tony Hoare called the null reference his ``billion-dollar mistake''
\citep{hoare2009null}.  In databases, null propagation is especially
insidious: \texttt{NULL + 1} evaluates to \texttt{NULL},
\texttt{NULL > 0} evaluates to \texttt{NULL} (not \texttt{false}), and
\texttt{NULL = NULL} evaluates to \texttt{NULL} (not \texttt{true}).
These semantics, while internally consistent, violate the expectations
of most programmers and are a leading source of silent data corruption.

\begin{example}
\begin{lstlisting}[language=VQL]
-- Schema: users.email is Nullable(String)
SELECT UPPER(DOCUMENT.email) FROM HEXAD ...
-- Level 4 error: UPPER expects String, got Nullable(String).
-- Fix: COALESCE(DOCUMENT.email, "unknown") or add WHERE ... IS NOT NULL

SELECT UPPER(COALESCE(DOCUMENT.email, "N/A")) FROM HEXAD ...
-- Accepted: COALESCE returns String (non-null guaranteed).
\end{lstlisting}
\end{example}

% ----------------------------------------------------------------------------
\subsection{Level 5: Injection-Proof Safety}\label{sec:L5}
% ----------------------------------------------------------------------------

\begin{definition}[\Safe{5}: Injection-Proof Safety]
A query $q$ is \Safe{5}-safe if and only if $q$ is \Safe{4}-safe and
there exists no function $f : \mathit{String} \to \mathit{Query}$ in
the language.  User-supplied values inhabit opaque parameter types
$\mathsf{Param}(\tau)$ that are structurally distinct from query syntax:
\[
  \Safe{5}(q) \iff \Safe{4}(q) \land
  \forall\, v \in \mathrm{params}(q).\; v : \mathsf{Param}(\tau_v) \land
  \neg \exists\, f.\; f : \mathit{String} \to \mathit{Query}
\]
\end{definition}

This is not sanitisation or escaping---it is a structural
impossibility.  The type system makes it impossible for user input to
be interpreted as query syntax, because the two inhabit disjoint types.

\paragraph{Mechanism.}
Parameters in \VQL{} are typed slots: \texttt{\$id :: UUID} declares a
parameter of type \texttt{UUID}.  The only way to supply a value is
through the typed parameter-binding API.  There is no string
interpolation operator and no \texttt{eval} function.  The parser is
the sole path from text to query, and it enforces all lower levels.

\paragraph{Failure mode without \LevelN{5}.}
Injection attacks---the persistent \#1 web application vulnerability
\citep{owasp2021top10}.  The 2017 Equifax breach exposed 147 million
records through an injection vector \citep{usgao2018equifax}.

\begin{example}
\begin{lstlisting}[language=VQL]
-- VQL-UT: injection impossible by construction
SELECT DOCUMENT.* FROM HEXAD ...
  WHERE DOCUMENT.id = $id :: UUID
-- $id is Param(UUID), not a string. Cannot contain VQL commands.
-- The type system prevents: $id = "'; DROP TABLE users; --"
\end{lstlisting}
\end{example}

% ----------------------------------------------------------------------------
\subsection{Level 6: Permission Safety}\label{sec:L6}
% ----------------------------------------------------------------------------

\begin{definition}[\Safe{6}: Permission Safety]
A query $q$ is \Safe{6}-safe if and only if $q$ is \Safe{5}-safe and
every field access and row access in $q$ is authorised by the
permission context $\Pi$:
\[
  \Safe{6}(q) \iff \Safe{5}(q) \land
  \forall\, a \in \mathrm{accesses}(q).\;
  \Pi \vdash a : \mathit{permitted}
\]
where $\Pi$ is a row-and-column-level access control policy and
$\mathrm{accesses}(q)$ is the set of data accesses in $q$.
\end{definition}

Permission safety ensures that the \emph{result type} of a query
reflects the caller's access rights.  A query executed by a user
without access to the \texttt{salary} column will have a result type
that does not include \texttt{salary}---the column is absent from the
type, not merely masked at runtime.

\paragraph{Mechanism.}
The \VQL{} type system parameterises queries by a permission context,
represented as a dependent type that encodes which columns and rows the
current principal may access.  The result type is computed as the
intersection of the selected fields and the permitted fields.

\begin{lstlisting}[language=Idris2]
data Permission : Schema -> Principal -> Type where
  Grant : (visible : SubSchema s) -> Permission s p

query : (perm : Permission schema principal) ->
        Query schema -> Result (restrict schema perm)
\end{lstlisting}

\paragraph{Failure mode without \LevelN{6}.}
Access control violations are discovered at runtime, after the query
has been planned and partially executed.  Or worse, overly permissive
defaults expose sensitive data.  The result type appears to contain
fields that the caller cannot actually read, leading to null values
or runtime exceptions.

% ----------------------------------------------------------------------------
\subsection{Level 7: Cardinality Safety}\label{sec:L7}
% ----------------------------------------------------------------------------

\begin{definition}[\Safe{7}: Cardinality Safety]
A query $q$ is \Safe{7}-safe if and only if $q$ is \Safe{6}-safe and
the expected result cardinality is part of the result type.  Let
$\kappa \in \{\mathsf{ExactlyOne}, \mathsf{AtMostOne},
\mathsf{NonEmpty}, \mathsf{Many}\}$ be the cardinality annotation.
Then:
\[
  \Safe{7}(q) \iff \Safe{6}(q) \land
  \exists\, \kappa.\; \tau(q) = \mathsf{Result}_\kappa(\sigma)
\]
where $\sigma$ is the row type and $\kappa$ is provable from schema
constraints and query structure.
\end{definition}

This is the first of the four novel safety levels.  Cardinality safety
encodes \emph{how many} results a query returns as part of the return
type, enabling the compiler to reject code that assumes a single result
when the query might return zero or many.

\paragraph{Mechanism.}
The cardinality is inferred from:
\begin{itemize}[nosep]
  \item \textbf{Unique constraints:} A filter on a unique key yields
    $\mathsf{ExactlyOne}$ (or $\mathsf{AtMostOne}$ if the filter is
    not a total function).
  \item \textbf{Aggregation:} \texttt{COUNT(*)} always returns
    $\mathsf{ExactlyOne}(\mathsf{Int})$.
  \item \textbf{LIMIT clauses:} \texttt{LIMIT 1} refines cardinality
    to $\mathsf{AtMostOne}$.
  \item \textbf{Explicit annotation:} \texttt{EXPECTING EXACTLY 1}
    declares the expected cardinality; the compiler verifies it.
\end{itemize}

\begin{lstlisting}[language=Idris2]
data Cardinality = ExactlyOne | AtMostOne | NonEmpty | Many

data Result : Cardinality -> Type -> Type where
  One  : a -> Result ExactlyOne a
  Opt  : Maybe a -> Result AtMostOne a
  NE   : (x : a) -> (xs : List a) -> Result NonEmpty a
  Bulk : List a -> Result Many a
\end{lstlisting}

\paragraph{Failure mode without \LevelN{7}.}
A login system queries users by email, assumes exactly one result, and
calls \texttt{.first()} on the result set.  When the query returns zero
rows, a null-pointer exception crashes the application.  When duplicate
emails exist (due to a missing unique constraint), the system
authenticates against the wrong account.

\begin{example}
\begin{lstlisting}[language=VQL]
SELECT DOCUMENT.* FROM HEXAD ...
  WHERE DOCUMENT.id = $id :: UUID
  EXPECTING EXACTLY 1
-- Return type: Result ExactlyOne UserRecord
-- Proved from: UNIQUE constraint on id

SELECT DOCUMENT.* FROM HEXAD ...
  WHERE DOCUMENT.email = $email :: String
  EXPECTING AT_MOST 1
-- Return type: Result AtMostOne UserRecord
-- Caller must handle the Nothing case.
\end{lstlisting}
\end{example}

% ----------------------------------------------------------------------------
\subsection{Level 8: Consistency Safety}\label{sec:L8}
% ----------------------------------------------------------------------------

\begin{definition}[\Safe{8}: Consistency Safety]
A query $q$ is \Safe{8}-safe if and only if $q$ is \Safe{7}-safe and
the effect set $E(q)$ is explicitly declared and enforced:
\[
  \Safe{8}(q) \iff \Safe{7}(q) \land
  E(q) \subseteq E_{\text{permitted}} \land
  \mathit{causal}(q)
\]
where $E(q) \subseteq \{\mathsf{Read}, \mathsf{Write}, \mathsf{DDL},
\mathsf{Network}\}$ is the effect set of $q$,
$E_{\text{permitted}}$ is the permitted effect set of the execution
context, and $\mathit{causal}(q)$ asserts that read-after-write
dependencies within $q$ are respected.
\end{definition}

Consistency safety unifies two concerns: \emph{effect tracking}
(distinguishing read from write operations at the type level) and
\emph{causal ordering} (ensuring that a write followed by a read within
the same logical operation returns the written value).

\paragraph{Mechanism.}
Effects are tracked using an algebraic effect system encoded in Idris\,2.
Each query is annotated with its effect set, and subsumption is checked:
a $\{\mathsf{Read}\}$ context cannot execute a $\mathsf{Write}$
operation.

\begin{lstlisting}[language=Idris2]
data Effect = Read | Write | DDL | Network

data Query : (effects : List Effect) -> Schema -> Type -> Type where
  Fetch  : Query [Read] s a
  Insert : Query [Read, Write] s a
  Alter  : Query [Read, Write, DDL] s a

-- Subsumption: a query with effects E can run in context E' iff E <= E'
subsumes : SubList e1 e2 -> Query e1 s a -> Query e2 s a
\end{lstlisting}

Causal ordering is tracked through a \emph{consistency model} type that
parameterises the connection.  A connection in \texttt{Causal} mode
guarantees read-your-writes; a connection in \texttt{Eventual} mode
does not.  The type of the result reflects the consistency level.

\paragraph{Failure mode without \LevelN{8}.}
Knight Capital Group lost \$440 million in 45 minutes on 1 August 2012
when a read-only monitoring function accidentally triggered write
operations on the New York Stock Exchange \citep{sec2013knight}.  The
root cause: nothing in the system prevented a ``read'' function from
performing writes.  Effect tracking makes this category of error
impossible.

\begin{example}
\begin{lstlisting}[language=VQL]
SELECT DOCUMENT.total FROM HEXAD ...
  EFFECTS { Read }
-- Accepted: SELECT is a read operation.

DELETE HEXAD 3f2504e0-...
  EFFECTS { Read }
-- Level 8 error: DELETE requires Write effect, but only Read is permitted.

DELETE HEXAD 3f2504e0-...
  EFFECTS { Read, Write }
-- Accepted: Write effect available.
\end{lstlisting}
\end{example}

% ----------------------------------------------------------------------------
\subsection{Level 9: Temporal Safety}\label{sec:L9}
% ----------------------------------------------------------------------------

\begin{definition}[\Safe{9}: Temporal Safety]
A query $q$ is \Safe{9}-safe if and only if $q$ is \Safe{8}-safe and
every temporal reference in $q$ carries a bound $\beta$ in the type:
\[
  \Safe{9}(q) \iff \Safe{8}(q) \land
  \forall\, t \in \mathrm{temporal}(q).\;
  \tau(t) = \mathsf{Temporal}(\sigma, \beta)
\]
where $\beta \in \{\mathsf{Freshness}(\delta), \mathsf{PointInTime}(t),
\mathsf{Range}(t_1, t_2), \mathsf{Snapshot}(\mathit{txn})\}$ is a
temporal bound.
\end{definition}

Temporal safety ensures that time-travel queries, freshness guarantees,
and bi-temporal constraints are encoded in the result type.  A query
result of type $\mathsf{Temporal}(\mathit{Price}, \mathsf{Freshness}(5s))$
carries a compile-time guarantee that the price data is at most 5
seconds old.

\paragraph{Mechanism.}
Temporal bounds are encoded as modal types indexed by temporal contexts:

\begin{lstlisting}[language=Idris2]
data TemporalBound : Type where
  Freshness  : (maxAge : Duration) -> TemporalBound
  PointInTime : (t : Timestamp) -> TemporalBound
  Range      : (from : Timestamp) -> (to : Timestamp) -> TemporalBound
  Snapshot   : (txn : TransactionId) -> TemporalBound

data Temporal : Type -> TemporalBound -> Type where
  MkTemporal : (value : a) -> (bound : TemporalBound) -> Temporal a bound
\end{lstlisting}

The runtime enforces the bound: if the most recent data for a
$\mathsf{Freshness}(\delta)$ query is older than $\delta$, the query
returns an error rather than stale data.

\paragraph{Failure mode without \LevelN{9}.}
The Flash Crash of 6 May 2010 saw the Dow Jones Industrial Average
drop nearly 1,000 points in minutes.  Contributing factors included
trading algorithms acting on stale price data as though it were current
\citep{kirilenko2017flash}.  Temporal safety would have encoded the
data's age in its type, making it impossible to use 5-minute-old data
in a context requiring sub-second freshness.

\paragraph{Bi-temporal queries.}
\VeriSimDB{} supports bi-temporal data: \emph{valid time} (when the
fact was true in the real world) and \emph{transaction time} (when it
was recorded in the database).  Temporal safety distinguishes these
at the type level, preventing accidental conflation.

\begin{example}
\begin{lstlisting}[language=VQL]
SELECT DOCUMENT.price FROM HEXAD ...
  WHERE DOCUMENT.ticker = "FTSE100"
  AS OF NOW - 5 SECONDS
  FRESH WITHIN 5s
-- Return type: Temporal(Price, Freshness(5s))
-- Guarantee: data is at most 5 seconds old.

SELECT TEMPORAL.* FROM HEXAD ...
  WHERE DOCUMENT.id = $id :: UUID
  BETWEEN "2026-01-01" AND "2026-01-31"
-- Return type: List(Temporal(Record, Range(2026-01-01, 2026-01-31)))
\end{lstlisting}
\end{example}

% ----------------------------------------------------------------------------
\subsection{Level 10: Provenance Safety}\label{sec:L10}
% ----------------------------------------------------------------------------

\begin{definition}[\Safe{10}: Provenance Safety]
A query $q$ is \Safe{10}-safe if and only if $q$ is \Safe{9}-safe and
every result $r$ of $q$ carries a provenance attestation $\pi$ that
is cryptographically verifiable:
\[
  \Safe{10}(q) \iff \Safe{9}(q) \land
  \forall\, r \in \mathrm{results}(q).\;
  \exists\, \pi.\; \mathit{verify}(\pi, r) = \mathsf{true}
\]
where $\pi$ is a cryptographic proof of origin (hash chain, Merkle
tree root, or zero-knowledge proof) and $\mathit{verify}$ is a
deterministic verification function.
\end{definition}

Provenance safety is the capstone of the decalogue.  It ensures that
query results are not merely correct but \emph{attestably}
correct---they carry cryptographic proof of their origin, their
transformation history, and their chain of custody.

\paragraph{Mechanism.}
Provenance is encoded as a type-level wrapper around query results:

\begin{lstlisting}[language=Idris2]
data ProvenanceAttestation : Type where
  HashChain : (root : Hash) -> (path : List Hash) -> ProvenanceAttestation
  MerkleProof : (root : Hash) -> (siblings : Vect n Hash) -> ProvenanceAttestation
  ZKProof : (proof : Bytes) -> (verifier : ContractId) -> ProvenanceAttestation

data Proved : Type -> Type where
  MkProved : (value : a) ->
             (attestation : ProvenanceAttestation) ->
             (valid : verify attestation value = True) ->
             Proved a
\end{lstlisting}

The critical feature is the \texttt{valid} field: it is a
\emph{proof term}---a value whose type asserts that the verification
succeeds.  This is not a runtime check; it is a compile-time
guarantee that the attestation is structurally valid.

\paragraph{Failure mode without \LevelN{10}.}
Without provenance safety, query results are opaque: the caller
cannot distinguish a genuine result from a tampered one.  In
regulatory contexts (financial auditing, medical records, legal
discovery), the inability to prove data provenance can have
legal and financial consequences.  In multi-party database
federations, provenance is essential for trust.

\paragraph{Integration with VeriSimDB.}
\VeriSimDB{}'s provenance modality natively records the
lineage of every data item---who created it, what transformations
it underwent, and what queries produced it.  Level~10 lifts this
runtime capability into the type system: the compiler ensures that
provenance chains are maintained through query compositions and that
results from the provenance modality carry verifiable attestations.

\begin{example}
\begin{lstlisting}[language=VQL]
SELECT DOCUMENT.*, PROVENANCE.chain
  FROM HEXAD 3f2504e0-...
  PROOF PROVENANCE(audit_trail)
  EXPECTING EXACTLY 1
  EFFECTS { Read, Proof }
-- Return type: Proved(Result ExactlyOne AuditRecord)
-- The result carries a cryptographic attestation of its origin.
-- The caller can verify the attestation without trusting the server.
\end{lstlisting}
\end{example}

% ============================================================================
\section{Cumulative Safety Ordering}\label{sec:lattice}
% ============================================================================

The ten safety levels form a \emph{totally ordered set} under
the subsumption relation.  We now formalise this structure and
prove the key properties.

\begin{definition}[Safety Ordering]
Define the binary relation $\Leq$ on safety levels by:
\[
  \LevelN{j} \Leq \LevelN{k} \iff j \leq k
\]
The pair $(\{\LevelN{1}, \ldots, \LevelN{10}\}, \Leq)$ is a
totally ordered set (a chain) with bottom element $\LevelN{1}$
and top element $\LevelN{10}$.
\end{definition}

\begin{theorem}[Subsumption]
For all queries $q$ and all safety levels $j < k$:
\[
  \Safe{k}(q) \implies \Safe{j}(q)
\]
\end{theorem}

\begin{proof}
By induction on $k$.  The base case $k = 1$ is vacuous (there is
no $j < 1$).  For the inductive step, observe that each definition
of $\Safe{k}$ includes $\Safe{k-1}(q)$ as a conjunct.  Thus
$\Safe{k}(q) \implies \Safe{k-1}(q)$, and by the inductive
hypothesis, $\Safe{k-1}(q) \implies \Safe{j}(q)$ for all $j < k-1$.
\end{proof}

\begin{corollary}[Maximum Safety]
$\Safe{10}(q) \implies \bigwedge_{j=1}^{9} \Safe{j}(q)$.
A Level~10-safe query is automatically safe at all lower levels.
\end{corollary}

\begin{proposition}[Strictness]
The ordering is strict: for each $k \in \{2, \ldots, 10\}$, there
exists a query $q_k$ such that $\Safe{k-1}(q_k) \land \neg
\Safe{k}(q_k)$.
\end{proposition}

\begin{proof}
We exhibit witnesses:
\begin{enumerate}[nosep,label=(\arabic*)]
  \item $q_2$: a syntactically valid query referencing a non-existent column.
  \item $q_3$: a schema-bound query comparing a string to an integer.
  \item $q_4$: a type-correct query applying a function to a nullable
    value without null handling.
  \item $q_5$: a null-safe query constructed via string concatenation.
  \item $q_6$: an injection-proof query accessing a column the caller lacks permission for.
  \item $q_7$: a permission-safe query where \texttt{.first()} is
    called on a potentially empty result.
  \item $q_8$: a cardinality-safe query that performs a write in a
    read-only context.
  \item $q_9$: a consistency-safe query that uses stale data without
    temporal bounds.
  \item $q_{10}$: a temporally safe query whose results cannot be
    verified for provenance.
\end{enumerate}
Each $q_k$ satisfies $\Safe{k-1}$ but violates $\Safe{k}$.
\end{proof}

\paragraph{Lattice extension.}
While the safety levels form a total order, the \emph{combinations of
safety properties that a system supports} form a bounded lattice.  A
system that supports Levels~1, 2, and~5 (but not~3 or~4) is below
$\LevelN{5}$ in our ordering---it has ``gaps.''  We define the
\emph{safety score} of a system as the highest $k$ such that Levels~1
through $k$ are all satisfied:

\begin{definition}[Safety Score]
The \emph{safety score} of a query language $L$ is:
\[
  \mathit{score}(L) = \max\{k \mid \forall\, j \leq k.\; L \text{ enforces } \Safe{j}\}
\]
\end{definition}

This definition penalises gaps: a system with Levels~1, 2, 3, and~5
but not~4 has a safety score of~3, because the chain is broken at~4.

% ============================================================================
\section{Type System Design}\label{sec:typesystem}
% ============================================================================

\VQL{}'s type system is built on four pillars, each contributing to
specific safety levels:

\begin{center}
\begin{tabular}{lll}
\toprule
\textbf{Pillar} & \textbf{Theoretical Foundation} & \textbf{Levels} \\
\midrule
Dependent types & Martin-L\"of Type Theory & L1--L6 \\
Refinement types & Liquid types \citep{rondon2008liquid} & L7 \\
Algebraic effects & Eff \citep{bauer2015programming} & L8 \\
Quantitative types & QTT \citep{atkey2018syntax,mcbride2016got} & L9--L10 \\
\bottomrule
\end{tabular}
\end{center}

\subsection{Dependent Types for Schema Binding (L1--L6)}

The schema is represented as a value-level term that is promoted to
the type level through Idris\,2's dependent types.  A query is a
term whose type depends on the schema:

\begin{lstlisting}[language=Idris2]
-- A query is parameterised by schema, permission, effects, and result type
data Query : (s : Schema) -> (p : Permission s) -> (e : Effects) ->
             (k : Cardinality) -> (t : Type) -> Type

-- The result type is computed from the query structure
typeOf : Query s p e k t -> t
\end{lstlisting}

Schema binding is a dependent function application: the column name
is a value, and the column type is computed from the schema.  If the
column does not exist, the dependent lookup fails at type-checking
time.

\subsection{Refinement Types for Cardinality (L7)}

Cardinality is expressed as a refinement on the result type.
The cardinality annotation is \emph{proved} from schema constraints
rather than merely declared:

\begin{lstlisting}[language=Idris2]
-- Prove that a filter on a unique key yields ExactlyOne
uniqueFilter : IsUnique s col -> Filter s col v ->
               Cardinality -> ExactlyOne
uniqueFilter uniq filt = ...  -- proof term

-- Prove that LIMIT 1 refines to AtMostOne
limitOne : HasLimit q 1 -> Cardinality -> AtMostOne
limitOne lim = ...  -- proof term
\end{lstlisting}

\subsection{Algebraic Effects for Consistency (L8)}

Effects form a join-semilattice:
\[
  \mathsf{Read} \Leq \mathsf{Read} \sqcup \mathsf{Write}
  \Leq \mathsf{Read} \sqcup \mathsf{Write} \sqcup \mathsf{DDL}
\]
A query with effect set $E$ can execute in a context with permitted
effects $E'$ if and only if $E \subseteq E'$.  This is checked at
compile time via a subsumption proof:

\begin{lstlisting}[language=Idris2]
data SubEffects : List Effect -> List Effect -> Type where
  SubNil  : SubEffects [] es
  SubCons : Elem e es -> SubEffects rest es -> SubEffects (e :: rest) es

run : SubEffects (queryEffects q) contextEffects -> Query ... -> IO Result
\end{lstlisting}

\subsection{Quantitative Type Theory for Linearity and Temporality (L9--L10)}

Idris\,2's QTT assigns a \emph{quantity} to each variable binding:
$0$ (erased at runtime), $1$ (used exactly once), or $\omega$
(unrestricted).  We exploit this for:

\begin{itemize}[nosep]
  \item \textbf{Connection linearity:} A connection handle has quantity
    $1$, ensuring it is opened and closed exactly once.
  \item \textbf{Transaction linearity:} A transaction handle is
    consumed by \texttt{COMMIT} or \texttt{ROLLBACK}, preventing
    double-commit.
  \item \textbf{Provenance tokens:} A provenance attestation has
    quantity $1$, ensuring it is consumed by exactly one verification
    step and cannot be reused across unrelated queries.
\end{itemize}

\begin{lstlisting}[language=Idris2]
-- Connection with quantity 1: must be used exactly once
open  : (1 _ : IO ()) -> IO (1 _ : Connection)
close : (1 _ : Connection) -> IO ()

-- Session types for transaction protocol
data SessionState = Ready | InTransaction | Closed
begin    : (1 _ : Session Ready) -> Session InTransaction
commit   : (1 _ : Session InTransaction) -> Session Ready
rollback : (1 _ : Session InTransaction) -> Session Ready
\end{lstlisting}

\subsection{Compositionality}

A key property of the type system is that safety levels \emph{compose}.
If query $q_1$ is \Safe{k}-safe and query $q_2$ is \Safe{k}-safe,
then their composition (join, union, subquery) is also \Safe{k}-safe,
provided the composition respects the invariants of each level.

\begin{theorem}[Composition Preservation]
For all $k \in \{1, \ldots, 10\}$ and composable queries $q_1$, $q_2$:
\[
  \Safe{k}(q_1) \land \Safe{k}(q_2) \land \mathit{composable}(q_1, q_2)
  \implies \Safe{k}(q_1 \bowtie q_2)
\]
where $\bowtie$ is any standard relational composition (join, union,
intersection, subquery).
\end{theorem}

\begin{proof}[Proof sketch]
By structural induction on the composition.  Each safety level's
predicate is closed under the standard relational algebra operations,
given that the type system propagates the relevant invariants.
Null safety propagates through joins (left join lifts right-hand
columns to nullable).  Cardinality composes via the product rule
($\mathsf{ExactlyOne} \bowtie \mathsf{ExactlyOne} = \mathsf{ExactlyOne}$
for key-foreign-key joins).  Effects compose via union.
Temporal bounds compose via intersection of time ranges.
Provenance attestations compose via Merkle tree concatenation.
Full proofs are provided in the supplementary material.
\end{proof}

% ============================================================================
\section{Integration with Multi-Modal Databases}\label{sec:multimodal}
% ============================================================================

\VeriSimDB{} stores data across eight modalities (the ``Octad''):
document, graph, vector, tensor, semantic, temporal, provenance, and
spatial.  Each modality has distinct data types, query semantics, and
consistency models.  The ten safety levels must therefore be applied
\emph{per modality} and \emph{across modalities}.

\subsection{The Octad Model}

\begin{center}
\begin{tabular}{lll}
\toprule
\textbf{Modality} & \textbf{Data Type} & \textbf{Query Semantics} \\
\midrule
Document   & JSON/BSON documents & Field access, full-text search \\
Graph      & Nodes and edges     & Traversal, path patterns \\
Vector     & Dense embeddings    & Similarity search, $k$-NN \\
Tensor     & $n$-d arrays        & Slice, reshape, contraction \\
Semantic   & ZKP contracts       & Satisfaction, verification \\
Temporal   & Versioned snapshots & Time-travel, bi-temporal \\
Provenance & Lineage chains      & Chain traversal, verification \\
Spatial    & Geometric objects   & Range, nearest, containment \\
\bottomrule
\end{tabular}
\end{center}

\subsection{Cross-Modal Safety}

A query that spans modalities must satisfy safety levels across the
modal boundary.  For example, a query that joins a document result
with a vector similarity search must ensure:

\begin{itemize}[nosep]
  \item \textbf{L2:} Both the document field and the vector field
    exist in the schema.
  \item \textbf{L3:} The join predicate compares compatible types
    (e.g., a document ID with a vector's source ID).
  \item \textbf{L4:} If the vector modality is optional for some
    entities, the join result is nullable.
  \item \textbf{L7:} The cardinality of a $k$-NN vector search
    is $\mathsf{AtMost}(k)$, and this propagates through the join.
  \item \textbf{L8:} Vector similarity search is a $\mathsf{Read}$
    operation; indexing is $\mathsf{Write} \cup \mathsf{Index}$.
  \item \textbf{L9:} Temporal bounds must be consistent across
    modalities---querying a document ``as of Tuesday'' while joining
    with a vector ``as of today'' is a type error.
  \item \textbf{L10:} Cross-modal provenance chains must be
    concatenated coherently.
\end{itemize}

The type system achieves this by parameterising queries by a
\emph{modality context}---a record of which modalities are active
and what their types are:

\begin{lstlisting}[language=Idris2]
record ModalityContext where
  constructor MkCtx
  document   : Maybe DocumentSchema
  graph      : Maybe GraphSchema
  vector     : Maybe VectorSchema
  temporal   : Maybe TemporalSchema
  provenance : Maybe ProvenanceSchema
  -- ... (spatial, tensor, semantic)

-- A cross-modal query is parameterised by the active modalities
crossModalQuery : (ctx : ModalityContext) ->
                  {auto docPresent : IsJust (document ctx)} ->
                  {auto vecPresent : IsJust (vector ctx)} ->
                  Query ...
\end{lstlisting}

The \texttt{auto} keyword instructs Idris\,2 to search for a proof
that the required modalities are present.  If a modality is absent,
the proof search fails and the query is rejected at compile time.

\subsection{Drift Detection}

\VeriSimDB{} supports \emph{drift detection}: when the same logical
entity has representations in multiple modalities that have diverged,
the drift is quantified and reported.  The VQL-UT grammar includes
drift policies (\texttt{STRICT}, \texttt{REPAIR}, \texttt{TOLERATE},
\texttt{LATEST}) that are reflected in the query's type:

\begin{lstlisting}[language=VQL]
SELECT DOCUMENT.*, VECTOR.embedding
  FROM FEDERATION /universities/*
  WITH DRIFT STRICT
  EFFECTS { Read }
  FRESH WITHIN 30s
-- If any node has drifted beyond threshold, the query fails.
-- Return type: Temporal(Result Many (DocRecord, VecEmbedding), Freshness(30s))
\end{lstlisting}

% ============================================================================
\section{Implementation}\label{sec:implementation}
% ============================================================================

\subsection{Architecture}

The \VQL{} implementation comprises four layers:

\begin{center}
\begin{tikzpicture}[
  box/.style={draw, rounded corners, minimum width=6cm, minimum height=1cm,
              fill=blue!8, font=\small},
  arr/.style={-{Stealth[length=5pt]}, thick},
  lbl/.style={font=\footnotesize\itshape, text=gray}
]
  \node[box] (parser)  at (0, 4.5) {Parser (ReScript) --- Level 1};
  \node[box] (binder)  at (0, 3.0) {Schema Binder (Idris\,2) --- Level 2};
  \node[box] (checker) at (0, 1.5) {Type Checker (Idris\,2) --- Levels 3--6};
  \node[box] (analyser) at (0, 0) {Effect/Temporal/Provenance Analyser (Idris\,2) --- Levels 7--10};
  \node[box, fill=green!8] (codegen)  at (0, -1.5) {FFI Codegen (Zig) --- C-ABI Query Plan};

  \draw[arr] (parser) -- (binder);
  \draw[arr] (binder) -- (checker);
  \draw[arr] (checker) -- (analyser);
  \draw[arr] (analyser) -- (codegen);

  \node[lbl, right=0.3cm] at (parser.east) {Source $\to$ AST};
  \node[lbl, right=0.3cm] at (binder.east) {AST $\to$ Bound AST};
  \node[lbl, right=0.3cm] at (checker.east) {Bound AST $\to$ Typed AST};
  \node[lbl, right=0.3cm] at (analyser.east) {Typed AST $\to$ Verified AST};
  \node[lbl, right=0.3cm] at (codegen.east) {Verified AST $\to$ Query Plan};
\end{tikzpicture}
\end{center}

\subsection{Parser (ReScript)}

The parser is a recursive-descent parser written in ReScript, extending
\VeriSimDB{}'s existing VQL v3.0 parser with the type-safety clauses
introduced in \Cref{sec:decalogue}.  ReScript was chosen for its strong
type inference, pattern matching, and compilation to efficient
JavaScript (enabling browser-based query editors with client-side
Level~1 checking).

The parser produces a typed AST in which each node carries its source
location, enabling precise error messages with suggestions (e.g.,
``Did you mean \texttt{email}?'' for a typo \texttt{emial}).

\subsection{Planned Verification Core (Idris\,2)}

The intended verification core is planned in Idris\,2, chosen for three
reasons:
\begin{enumerate}[nosep]
  \item \textbf{Dependent types} enable schema-level type checking
    (Levels~2--6).
  \item \textbf{Totality checking} (\texttt{\%default total}) ensures
    the type checker itself always terminates.
  \item \textbf{Quantitative Type Theory} (QTT) natively supports
    linear resource tracking (Levels~9--10).
\end{enumerate}

The long-term verification core is intended to be structured as a series
of passes, each corresponding to one or more safety levels. The current
repository contains draft Idris\,2 encodings for parts of the model and
ABI, not a checked end-to-end Levels~2--10 implementation.

\subsection{Planned FFI Bridge (Zig)}

The intended architecture compiles a verified AST to a C-ABI query plan
using a code generator written in Zig.  Zig was chosen for its zero-overhead C interoperability,
cross-compilation support, and safety features (bounds checking,
no hidden control flow).  The query plan is a sequence of low-level
operations that the \VeriSimDB{} engine can execute directly, with
no runtime type checking overhead in the design target. The current
repository does not yet provide this full checked pipeline.

\subsection{Performance Considerations}

A natural concern is that ten levels of compile-time checking impose
unacceptable compilation overhead.  We observe that:

\begin{itemize}[nosep]
  \item Levels~1--3 add negligible overhead (parsing and type checking
    are linear or near-linear in query size).
  \item Level~4 (null tracking) adds overhead proportional to the
    number of joins (each join may lift columns to nullable).
  \item Levels~5--6 are structural (enforced by the type system's
    lack of unsafe coercions) and add zero checking overhead.
  \item Level~7 (cardinality) requires constraint propagation, which
    is bounded by the number of schema constraints.
  \item Level~8 (effects) is a simple set-inclusion check.
  \item Level~9 (temporal) adds bounds checking proportional to the
    number of temporal references.
  \item Level~10 (provenance) adds overhead proportional to the
    provenance chain length, but verification is deferred to a
    separate pass.
\end{itemize}

These are design expectations rather than established measurements.
The current repository does not yet provide the complete compiler path
needed to validate compile-time-only enforcement or zero-overhead code
generation across all ten levels.

% ============================================================================
\section{Comparison with Existing Query Languages}\label{sec:comparison}
% ============================================================================

We evaluate existing query languages against the decalogue, assigning
each a safety score (the highest level $k$ such that all levels
1 through $k$ are satisfied).

\begin{table}[ht]
\centering
\caption{Safety scores of existing query languages.  A checkmark
(\checkmark) indicates the level is fully satisfied; a half-mark
($\sim$) indicates partial satisfaction; a cross (\texttimes)
indicates the level is not addressed.  The safety score is the
highest unbroken prefix.}
\label{tab:comparison}
\small
\begin{tabular}{l*{10}{c}c}
\toprule
\textbf{Language} &
\rotatebox{70}{\LevelN{1} Parse} &
\rotatebox{70}{\LevelN{2} Schema} &
\rotatebox{70}{\LevelN{3} Types} &
\rotatebox{70}{\LevelN{4} Null} &
\rotatebox{70}{\LevelN{5} Injection} &
\rotatebox{70}{\LevelN{6} Permission} &
\rotatebox{70}{\LevelN{7} Cardinality} &
\rotatebox{70}{\LevelN{8} Consistency} &
\rotatebox{70}{\LevelN{9} Temporal} &
\rotatebox{70}{\LevelN{10} Provenance} &
\textbf{Score} \\
\midrule
Raw SQL strings    & \texttimes & \texttimes & \texttimes & \texttimes & \texttimes & \texttimes & \texttimes & \texttimes & \texttimes & \texttimes & 0 \\
Prepared stmts     & \checkmark & \texttimes & \texttimes & \texttimes & $\sim$     & \texttimes & \texttimes & \texttimes & \texttimes & \texttimes & 1 \\
Django ORM         & \checkmark & $\sim$     & $\sim$     & \texttimes & \checkmark & \texttimes & \texttimes & \texttimes & \texttimes & \texttimes & 1 \\
Prisma             & \checkmark & \checkmark & \checkmark & $\sim$     & \checkmark & $\sim$     & \texttimes & \texttimes & \texttimes & \texttimes & 3 \\
jOOQ               & \checkmark & \checkmark & \checkmark & $\sim$     & \checkmark & \texttimes & \texttimes & \texttimes & \texttimes & \texttimes & 3 \\
LINQ               & \checkmark & \checkmark & \checkmark & $\sim$     & \checkmark & \texttimes & \texttimes & \texttimes & \texttimes & \texttimes & 3 \\
Slick/Quill        & \checkmark & \checkmark & \checkmark & $\sim$     & \checkmark & \texttimes & \texttimes & \texttimes & \texttimes & \texttimes & 3 \\
GraphQL            & \checkmark & \checkmark & \checkmark & $\sim$     & \checkmark & $\sim$     & \texttimes & \texttimes & \texttimes & \texttimes & 3 \\
Ur/Web             & \checkmark & \checkmark & \checkmark & \checkmark & \checkmark & \checkmark & \texttimes & \texttimes & \texttimes & \texttimes & 6 \\
EdgeQL             & \checkmark & \checkmark & \checkmark & \checkmark & \checkmark & $\sim$     & $\sim$     & \texttimes & \texttimes & \texttimes & 5 \\
PRQL               & \checkmark & $\sim$     & $\sim$     & \texttimes & \checkmark & \texttimes & \texttimes & \texttimes & \texttimes & \texttimes & 1 \\
Malloy             & \checkmark & \checkmark & \checkmark & $\sim$     & \checkmark & \texttimes & \texttimes & \texttimes & \texttimes & \texttimes & 3 \\
\textbf{VQL-UT (design target)} & \checkmark & \checkmark & \checkmark & \checkmark & \checkmark & \checkmark & \checkmark & \checkmark & \checkmark & \checkmark & \textbf{10} \\
\bottomrule
\end{tabular}
\end{table}

\paragraph{Observations.}
\begin{enumerate}[nosep]
  \item No existing language achieves a safety score above~6.
    Ur/Web comes closest, but lacks cardinality, consistency,
    temporal, and provenance safety.
  \item Most production query systems achieve a score of~3:
    they provide parse safety, schema binding, and type checking,
    but lack null safety (or provide it only partially).
  \item Injection-proof safety (L5) is widely supported---but because
    the chain is broken at L4 (null safety), the effective safety
    score remains low.
  \item Levels~7--10 are entirely unaddressed in the surveyed
    systems, confirming their novelty.
\end{enumerate}

\paragraph{A note on fairness.}
Several of these systems (notably Prisma, EdgeQL, and Malloy) are
under active development and may address additional levels in future
versions.  The comparison reflects the state of each system at the
time of writing (March 2026).  Our intent is not to disparage these
excellent systems but to identify the frontier they have not yet
reached.

The final row describes \VQL{} as a design target rather than a current
repository achievement. The present repo does not yet implement or prove
all ten levels end to end.

% ============================================================================
\section{The Classification Insight}\label{sec:classification}
% ============================================================================

A central observation of this work is that database query bugs form
a \emph{total order by severity and subtlety}.  This is not a
coincidence---it reflects the structure of the query processing
pipeline:

\begin{enumerate}[nosep]
  \item A query must be parsed before it can be bound to a schema.
  \item A query must be schema-bound before its operations can be
    type-checked.
  \item A query must be type-checked before null propagation can be
    tracked.
  \item A query must be null-safe before injection-proofing is
    meaningful (injected SQL exploits type coercion and null confusion).
  \item A query must be injection-proof before permission checking
    is meaningful (an injected query bypasses all permissions).
  \item Permissions must be resolved before cardinality can be
    computed (the visible rows determine the result count).
  \item Cardinality must be known before effects can be tracked
    (a write that produces $\mathsf{ExactlyOne}$ result is different
    from a write that produces $\mathsf{Many}$).
  \item Effects must be tracked before temporal bounds are meaningful
    (a write invalidates previous temporal guarantees).
  \item Temporal bounds must be established before provenance can
    be attested (a provenance chain includes timestamps).
\end{enumerate}

This chain is not merely pedagogical---it reflects genuine
\emph{logical dependencies} between the safety properties.  Each
level's guarantee is \emph{vacuous} without the guarantees below it.
Cardinality safety (L7) is meaningless if the query might reference
non-existent columns (L2 failure), because the cardinality of a
malformed query is undefined.  Provenance safety (L10) is meaningless
if the temporal context is wrong (L9 failure), because the provenance
chain would attest to the wrong version of the data.

This observation---that the ten properties form a total order of
logical dependency, not merely a checklist---is the central
taxonomic contribution of this paper.

% ============================================================================
\section{Formal Properties}\label{sec:formal}
% ============================================================================

We state several target properties of the \VQL{} type system. Full
proofs are \emph{not} yet mechanised in the current repository; the
statements below should be read as the intended metatheory for the
proposed system rather than as checked results.

\begin{theorem}[Type Soundness]
If a \VQL{} query $q$ type-checks at safety level $k$, then
execution of $q$ will not exhibit any failure mode associated with
levels $1$ through $k$.
\end{theorem}

\begin{proof}[Proof sketch]
By a standard progress-and-preservation argument adapted to the query
domain.  Progress: a well-typed query can always take a step (execute
a query plan instruction) or has completed.  Preservation: each step
maintains all safety invariants.  The proof is stratified by level:
each level's invariants are preserved independently, and the
subsumption theorem (\Cref{sec:lattice}) ensures that lower-level
invariants are maintained by higher-level checking.
\end{proof}

\begin{theorem}[Decidability]
For each safety level $k \in \{1, \ldots, 10\}$, the predicate
$\Safe{k}(q)$ is decidable for finite schemas and finite queries.
\end{theorem}

\begin{proof}[Proof sketch]
Level~1 is decidable because the grammar is context-free (CYK or
Earley parsing).  Levels~2--6 are decidable because the schema is
finite and the type system is stratified (no recursive types at the
schema level).  Level~7 is decidable because constraint propagation
over finite schemas terminates.  Level~8 is decidable because effect
inclusion on finite effect sets is decidable.  Level~9 is decidable
because temporal bound comparison is decidable over a discrete time
domain.  Level~10 is decidable because cryptographic verification
is a deterministic computation on finite inputs.
\end{proof}

\begin{theorem}[Zero Runtime Overhead]
A \Safe{10}-safe query compiles to a query plan that contains no
runtime safety checks.  The generated plan is identical in structure
and performance to a plan generated from an unchecked query.
\end{theorem}

\begin{proof}[Proof sketch]
All safety checks are erased during compilation (Idris\,2's erasure
analysis removes all proof terms with quantity~$0$).  The Zig code
generator receives a verified AST and produces a query plan containing
only data operations.  The safety invariants have been \emph{proved},
not \emph{checked}---they need not be verified again at runtime.
\end{proof}

% ============================================================================
\section{Future Work}\label{sec:future}
% ============================================================================

\subsection{Automated Level Inference}

Currently, \VQL{} queries may include explicit annotations for
higher safety levels (e.g., \texttt{EXPECTING EXACTLY 1},
\texttt{EFFECTS \{Read\}}).  We plan to implement \emph{automatic
level inference}: the compiler determines the highest safety level
a query satisfies without explicit annotations, and reports this as a
``safety score'' alongside the type-checking result.  This would
enable gradual adoption: existing VQL queries are assigned scores
without modification, and developers incrementally add annotations to
raise their scores.

\subsection{Safety Certificate Generation}

We envision \VQL{} producing \emph{safety certificates}---machine-readable
proofs that a query satisfies a given safety level.  These certificates
could be:
\begin{itemize}[nosep]
  \item Attached to query plans for auditing (``this plan was proved
    \Safe{8} at compilation time'').
  \item Verified by the database engine before execution (defence in
    depth).
  \item Stored alongside query results for provenance (extending
    Level~10 to the query itself, not just the data).
\end{itemize}

\subsection{Cross-Language Safety Propagation}

When a \VQL{} query is embedded in a host language (ReScript, Elixir,
Rust), the safety guarantees should propagate to the host.  For
example, a Level~7 (\texttt{ExactlyOne}) result should be represented
as a non-optional type in ReScript, not as a list that the caller
must destructure.  We plan to implement typed code generation for
multiple host languages, preserving the safety level through the FFI
boundary.

\subsection{Empirical Evaluation}

While we have provided formal proofs and a reference implementation
strategy, a comprehensive empirical evaluation remains future work.
We plan to measure:
\begin{itemize}[nosep]
  \item Compilation time overhead at each safety level.
  \item Runtime performance of generated query plans versus hand-written SQL.
  \item Developer experience: time to diagnose and fix bugs with and
    without safety levels.
  \item Bug detection rate: how many real-world database bugs would be
    caught at each level.
\end{itemize}

\subsection{Extension to Distributed Queries}

\VeriSimDB{} supports federated queries across multiple nodes.
Distributed queries introduce additional failure modes: network
partitions, partial results, and replica divergence.  We conjecture
that at least two additional safety levels may be needed:
$\LevelN{11}$~(partition-tolerance safety) and
$\LevelN{12}$~(convergence safety).  Whether these can be encoded
as type-level invariants without sacrificing decidability is an
open question.

% ============================================================================
\section{Conclusion}\label{sec:conclusion}
% ============================================================================

We have presented the \emph{Decalogue of Database Query Safety}:
ten progressive safety levels that form a total order under
subsumption.  The first six levels---parse-time safety, schema-binding
safety, type-compatible operations, null safety, injection-proof
safety, and permission safety---are individually addressed in the
existing literature, though no prior system unifies all six.  The
last four levels---cardinality safety, consistency safety, temporal
safety, and provenance safety---are, to our knowledge, the first to
be formalised as type-level invariants in a query language.

The key insight is that database bugs are not a cloud of independent
failure modes but a \emph{chain}: each class of bug presupposes the
absence of all less severe classes.  This total ordering enables a
meaningful ``safety score'' for any query language---the highest level
at which the chain is unbroken. Existing deployed query languages
appear to score between 0 and~6 in the comparison above, while \VQL{}
is proposed as a design target for score~10 rather than a present
repository achievement.

The practical consequence, if the full system is completed, would be a
query language in which these guarantees are enforced at compile time
with minimal or zero runtime overhead. The current repository does not
yet justify claiming that end-to-end state.

The theoretical consequence is a framework for classifying and
comparing query language safety properties.  The decalogue provides a
common vocabulary for discussing what a query language guarantees, what
it leaves to convention, and what it cannot address.  We hope this
taxonomy proves useful beyond the specific context of \VQL{} and
\VeriSimDB{}, as a tool for reasoning about the safety of any system
that transforms data.

% ============================================================================
% References
% ============================================================================
\begin{thebibliography}{30}

\bibitem[Atkey(2018)]{atkey2018syntax}
Atkey, R. (2018).
\newblock Syntax and semantics of quantitative type theory.
\newblock In \emph{Proceedings of the 33rd Annual ACM/IEEE Symposium on Logic
  in Computer Science (LICS)}, pp.~56--65.

\bibitem[Bauer and Pretnar(2015)]{bauer2015programming}
Bauer, A. and Pretnar, M. (2015).
\newblock Programming with algebraic effects and handlers.
\newblock \emph{Journal of Logical and Algebraic Methods in Programming},
  84(1):108--123.

\bibitem[Bauer et~al.(2015)]{bauer2015java}
Bauer, C., King, G., and Gregory, G. (2015).
\newblock \emph{Java Persistence with Hibernate}.
\newblock Manning Publications, 2nd edition.

\bibitem[Brady(2013)]{brady2013idris}
Brady, E. (2013).
\newblock Idris, a general-purpose dependently typed programming language:
  Design and implementation.
\newblock \emph{Journal of Functional Programming}, 23(5):552--593.

\bibitem[Buneman et~al.(2001)]{buneman2001and}
Buneman, P., Khanna, S., and Tan, W.-C. (2001).
\newblock Why and where: A characterization of data provenance.
\newblock In \emph{International Conference on Database Theory (ICDT)},
  pp.~316--330.

\bibitem[Cheney et~al.(2009)]{cheney2009provenance}
Cheney, J., Chiticariu, L., and Tan, W.-C. (2009).
\newblock Provenance in databases: Why, how, and where.
\newblock \emph{Foundations and Trends in Databases}, 1(4):379--474.

\bibitem[Chlipala(2010)]{chlipala2010ur}
Chlipala, A. (2010).
\newblock Ur: Statically-typed metaprogramming with type-level record
  computation.
\newblock In \emph{ACM SIGPLAN Conference on Programming Language Design and
  Implementation (PLDI)}, pp.~122--133.

\bibitem[Cockx et~al.(2021)]{cockx2021idris2}
Cockx, J., Brady, E., and Atkey, R. (2021).
\newblock Idris 2: Quantitative type theory in practice.
\newblock In \emph{European Conference on Object-Oriented Programming
  (ECOOP)}, pp.~9:1--9:26.

\bibitem[Curino et~al.(2008)]{curino2008automating}
Curino, C., Moon, H.~J., and Zaniolo, C. (2008).
\newblock Automating the database schema evolution process.
\newblock \emph{The VLDB Journal}, 22(1):73--98.

\bibitem[EdgeDB(2023)]{edgedb2023docs}
EdgeDB (2023).
\newblock EdgeQL documentation.
\newblock \url{https://www.edgedb.com/docs/edgeql/}.

\bibitem[GraphQL Foundation(2023)]{graphql2023spec}
GraphQL Foundation (2023).
\newblock GraphQL specification.
\newblock \url{https://spec.graphql.org/}.

\bibitem[Green et~al.(2007)]{green2007provenance}
Green, T.~J., Karvounarakis, G., and Tannen, V. (2007).
\newblock Provenance semirings.
\newblock In \emph{ACM SIGMOD-SIGACT-SIGART Symposium on Principles of Database
  Systems (PODS)}, pp.~31--40.

\bibitem[Halfond et~al.(2006)]{halfond2006classification}
Halfond, W.~G., Viegas, J., and Orso, A. (2006).
\newblock A classification of {SQL}-injection attacks and countermeasures.
\newblock In \emph{IEEE International Symposium on Secure Software Engineering
  (ISSSE)}, pp.~13--15.

\bibitem[Hoare(2009)]{hoare2009null}
Hoare, C. A.~R. (2009).
\newblock Null references: The billion dollar mistake.
\newblock Presentation at QCon London 2009.

\bibitem[Holovaty and Kaplan-Moss(2009)]{holovaty2009django}
Holovaty, A. and Kaplan-Moss, J. (2009).
\newblock \emph{The Definitive Guide to Django}.
\newblock Apress, 2nd edition.

\bibitem[Ioffe(2017)]{ioffe2017quill}
Ioffe, F. (2017).
\newblock Quill: Compile-time language integrated queries for Scala.
\newblock \url{https://getquill.io/}.

\bibitem[jOOQ(2023)]{jooq2023docs}
jOOQ (2023).
\newblock jOOQ documentation.
\newblock \url{https://www.jooq.org/doc/latest/manual/}.

\bibitem[Kim et~al.(2013)]{kim2013silent}
Kim, W., Choi, B.-J., Hong, E.-K., Kim, S.-K., and Lee, D. (2013).
\newblock A taxonomy of dirty data.
\newblock \emph{Data Mining and Knowledge Discovery}, 7(1):81--99.

\bibitem[Kirilenko et~al.(2017)]{kirilenko2017flash}
Kirilenko, A., Kyle, A.~S., Samadi, M., and Tuzun, T. (2017).
\newblock The flash crash: High-frequency trading in an electronic market.
\newblock \emph{The Journal of Finance}, 72(3):967--998.

\bibitem[Leveson and Turner(1993)]{leveson1993therac}
Leveson, N.~G. and Turner, C.~S. (1993).
\newblock An investigation of the {Therac-25} accidents.
\newblock \emph{IEEE Computer}, 26(7):18--41.

\bibitem[Malloy(2023)]{malloy2023docs}
Malloy (2023).
\newblock Malloy documentation.
\newblock \url{https://www.malloydata.dev/documentation}.

\bibitem[McBride(2016)]{mcbride2016got}
McBride, C. (2016).
\newblock I got plenty o' nuttin'.
\newblock In \emph{A List of Successes That Can Change the World}, pp.~207--233.

\bibitem[Meijer et~al.(2006)]{meijer2006linq}
Meijer, E., Beckman, B., and Bierman, G. (2006).
\newblock {LINQ}: Reconciling object, relations and {XML} in the {.NET}
  framework.
\newblock In \emph{ACM SIGMOD International Conference on Management of Data},
  pp.~706--706.

\bibitem[Munroe(2007)]{xkcd327}
Munroe, R. (2007).
\newblock Exploits of a Mom (xkcd \#327).
\newblock \url{https://xkcd.com/327/}.

\bibitem[Oury and Swierstra(2008)]{oury2008power}
Oury, N. and Swierstra, W. (2008).
\newblock The power of {Pi}.
\newblock In \emph{ACM SIGPLAN International Conference on Functional
  Programming (ICFP)}, pp.~39--50.

\bibitem[OWASP(2021)]{owasp2021top10}
OWASP Foundation (2021).
\newblock {OWASP} Top Ten 2021.
\newblock \url{https://owasp.org/www-project-top-ten/}.

\bibitem[PRQL(2023)]{prql2023docs}
PRQL (2023).
\newblock PRQL language documentation.
\newblock \url{https://prql-lang.org/}.

\bibitem[Prisma(2023)]{prisma2023docs}
Prisma (2023).
\newblock Prisma documentation.
\newblock \url{https://www.prisma.io/docs/}.

\bibitem[Rain~Forest~Puppy(1998)]{rfp1998nt}
Rain~Forest~Puppy (1998).
\newblock {NT} web technology vulnerabilities.
\newblock \emph{Phrack Magazine}, 8(54).

\bibitem[Rondon et~al.(2008)]{rondon2008liquid}
Rondon, P.~M., Kawaguchi, M., and Jhala, R. (2008).
\newblock Liquid types.
\newblock In \emph{ACM SIGPLAN Conference on Programming Language Design and
  Implementation (PLDI)}, pp.~159--169.

\bibitem[SEC(2013)]{sec2013knight}
U.S.~Securities and Exchange Commission (2013).
\newblock In the matter of Knight Capital Americas {LLC}: Administrative
  proceeding file no.~3-15570.

\bibitem[Slick(2023)]{slick2023docs}
Slick (2023).
\newblock Slick: Functional relational mapping for Scala.
\newblock \url{https://scala-slick.org/doc/3.5.0/}.

\bibitem[U.S.~GAO(2018)]{usgao2018equifax}
U.S.~Government Accountability Office (2018).
\newblock Data protection: Actions taken by Equifax and federal agencies in
  response to the 2017 breach.
\newblock Report GAO-18-559.

\bibitem[Weirich et~al.(2017)]{weirich2017dependent}
Weirich, S., Voizard, A., de~Amorim, P. H.~A., and Eisenberg, R.~A. (2017).
\newblock A specification for dependent types in Haskell.
\newblock In \emph{ACM SIGPLAN International Conference on Functional
  Programming (ICFP)}, pp.~31:1--31:29.

\end{thebibliography}

% ============================================================================
% Appendix
% ============================================================================
\appendix

\section{VQL-UT Grammar (EBNF Extract)}\label{app:grammar}

The following is an abbreviated extract of the VQL-UT grammar,
showing the top-level structure and the type-safety clauses.  The
complete grammar (583 lines of EBNF) is available in the project
repository.

\begin{lstlisting}[language={},basicstyle=\footnotesize\ttfamily,
  backgroundcolor=\color{bglight},frame=single]
statement = query | mutation ;

query = select_clause,
        from_clause,
        [typed_by_clause],       (* Level 2: Schema-binding *)
        [where_clause],
        [group_by_clause],
        [having_clause],
        [proof_clause],
        [order_by_clause],
        [limit_clause],
        [offset_clause],
        [returns_clause],        (* Level 6: Result-type *)
        [cardinality_clause],    (* Level 7: Cardinality *)
        [effects_clause],        (* Level 8: Effect-tracking *)
        [freshness_clause],      (* Level 9: Temporal safety *)
        [linearity_clause] ;     (* Level 10: Linearity *)

typed_by_clause = 'TYPED', 'BY', schema_ref ;
returns_clause  = 'RETURNS', '(', returns_column_list, ')' ;
cardinality_clause = 'EXPECTING', cardinality_spec ;
effects_clause  = 'EFFECTS', '{', effect_list, '}' ;
freshness_clause = 'FRESH', 'WITHIN', duration ;
linearity_clause = linearity_consume | linearity_session
                 | linearity_usage ;
\end{lstlisting}

\section{Summary of Notation}\label{app:notation}

\begin{center}
\begin{tabular}{ll}
\toprule
\textbf{Symbol} & \textbf{Meaning} \\
\midrule
$\Safe{k}(q)$ & Query $q$ satisfies safety level $k$ \\
$\LevelN{k}$ & Safety level $k$ \\
$\Leq$ & Subsumption ordering on safety levels \\
$\Sigma$ & Database schema \\
$\Pi$ & Permission context \\
$E(q)$ & Effect set of query $q$ \\
$\kappa$ & Cardinality annotation \\
$\beta$ & Temporal bound \\
$\pi$ & Provenance attestation \\
$\tau(e)$ & Type of expression $e$ \\
$\mathrm{refs}(q)$ & Schema references in query $q$ \\
$\mathrm{dom}(\Sigma)$ & Domain of schema $\Sigma$ \\
$\mathcal{L}(G)$ & Language generated by grammar $G$ \\
$\bowtie$ & Relational composition (join/union/subquery) \\
\bottomrule
\end{tabular}
\end{center}

\end{document}
